\documentclass{article}

% neurips_2023.sty is present in this directory.
% Anonymous mode is the default and is REQUIRED for BayLearn (double-blind).
\usepackage{neurips_2023}

\usepackage[utf8]{inputenc}
\usepackage[T1]{fontenc}
\usepackage{hyperref}
\usepackage{url}
\usepackage{booktabs}
\usepackage{amsfonts}
\usepackage{nicefrac}
\usepackage{microtype}

% Blank the PDF metadata so no author name is written into the file properties.
\hypersetup{pdfauthor={}, pdfsubject={}, pdfkeywords={}, pdfcreator={}, pdfproducer={}}

\title{Internal Validation Does Not Establish Support for\\
Individual Predictions: A 26-Case Diagnostic Study}

% Leave author block anonymous. Do not un-anonymize for submission.
\author{%
  Anonymous Author(s) \\
  Affiliation \\
  Address \\
  \texttt{email} \\
}

\begin{document}

\maketitle

\begin{abstract}
High aggregate discrimination and stable internal validation performance do not show
that a model's estimate is supported for every patient. This matters because
vasodilators used for idiopathic pulmonary arterial hypertension (IPAH) can cause
life-threatening pulmonary edema when pulmonary veno-occlusive disease (PVOD) is
misdiagnosed as IPAH. We reanalyzed published data from 166 patients (26 with PVOD and
140 with idiopathic or heritable PAH). Rather than apply one GDF15 cutoff, we estimated
patient-specific likelihood ratios and withheld point estimates when local data lacked
observations from both groups. GDF15 discriminated PVOD better than any routine clinical
variable (AUC, 0.956 versus 0.673). Yet values near the published cutoff provided little
evidence, and observations far from the cutoff accounted for much of the aggregate
performance. Internal validation changed the aggregate score by 0.0001, whereas
leave-one-case-out refitting changed the estimate for the most extreme patient by 15\%.
Aggregate generalization and local support are separate properties. Diagnostic-model
reports should include per-patient influence and abstention when support is
insufficient.
\end{abstract}

\section{Setting}

PVOD and IPAH have similar clinical presentations, but their treatment differs. PAH
vasodilators can cause life-threatening pulmonary edema in PVOD. A recent research
letter~\cite{source} reported that GDF15 distinguishes the two conditions and proposed a
threshold of 1{,}658\,pg/mL: values above the threshold suggest PVOD, and values below it
suggest IPAH.

Among 322 published patients, 166 met the criteria for this analysis: 26 had confirmed
PVOD and 140 had idiopathic or heritable PAH. A single threshold discards information
about distance from the cutoff. We instead estimated a likelihood ratio: how often a
value near the patient's measurement occurs in PVOD relative to the comparison group. We
estimated the ratio three ways --- a two-parameter lognormal fitted separately to each
group with unequal variances, a model-free banded estimate with a continuity correction,
and logistic regression on log GDF15 with a quadratic curvature term --- and obtained
intervals by stratified bootstrap (4{,}000 resamples, resampling within each group). We
then assessed whether both groups had observations near the queried value. When local
support was absent, we reported a conservative bound rather than a point estimate.

\section{Results}

\textbf{GDF15 had the highest discrimination.} We evaluated each routinely collected
clinical variable separately (Table~\ref{tab:auc}). The best, six-minute walk distance,
had an AUC of 0.673; GDF15 had an AUC of 0.956 (95\% CI, 0.912--0.990). GDF15 had little
correlation with hemodynamic and blood measurements, suggesting that it provides
additional information. We used complete cases without imputation, so samples differed
across variables.

\begin{table}[t]
\caption{AUC for distinguishing PVOD from idiopathic or heritable PAH.}
\label{tab:auc}
\centering
\small
\begin{tabular}{lrr}
\toprule
Variable & Patients & Separation (AUC) \\
\midrule
GDF15 (blood biomarker)        & 166 & 0.956 \\
Six-minute walk distance       & 112 & 0.673 \\
Mean pulmonary artery pressure & 157 & 0.617 \\
Pulmonary vascular resistance  & 142 & 0.602 \\
Cardiac index                  &  82 & 0.567 \\
NT-proBNP                      &  65 & 0.558 \\
Patient age                    & 166 & 0.522 \\
\bottomrule
\end{tabular}
\end{table}

\textbf{Values near the cutoff provided little evidence.} The published threshold had
81\% sensitivity and 98\% specificity, but most PVOD values were far above
1{,}658\,pg/mL. Patient-specific analysis identified an inconclusive range of
approximately 1{,}000 to 2{,}500\,pg/mL. At exactly 1{,}658\,pg/mL, a pretest probability
of 20\% increased to 37\%; the estimated interval was 19\% to 46\% and included the
pretest probability. Thus, uncertainty spanned a range rather than a single cutoff.
Values below 500\,pg/mL provided the strongest evidence against PVOD in this dataset, a
result not emphasized in the original report.

\textbf{Individual cases affected estimates at the upper tail.} Internal validation by
0.632+ bootstrap (4{,}000 resamples) changed the aggregate score by 0.0001 but did not
measure how strongly an individual case supports a patient-specific estimate. We refit
the model 26 times, omitting one PVOD case at a time. Omitting the patient with the
highest GDF15 value reduced that patient's estimate by 15\%; the influence of omitted
cases decreased rapidly at lower values.

\textbf{Comparison patients were absent at high GDF15 values.} Eighteen of the 26
patients with PVOD had values above 2{,}500\,pg/mL, but the comparison group had none
above approximately 2{,}000\,pg/mL. The likelihood ratio at high values could therefore
be estimated only as a lower bound because positivity was limited. We estimated that
values in this range were at least 33 times more common in PVOD than in the comparison
group. The data support strong discrimination but do not identify an upper limit for the
likelihood ratio.

\section{Discussion}

\textbf{Interpretation.} A lower bound of 33 may understate the evidence for PVOD if the
true likelihood ratio is larger. However, a point estimate without empirical support
conveys false precision. Both directions of diagnostic error have serious consequences:
underdiagnosis of PVOD can lead to harmful vasodilator treatment, whereas overdiagnosis
can withhold effective treatment from a patient with IPAH. This analysis does not
determine the appropriate clinical action threshold.

\textbf{Implications.} Good aggregate performance does not establish support for every
patient-level prediction, especially near the edge of the observed data. The problem
applies whenever a model is evaluated in aggregate but used one patient at a time.
Per-patient influence and local support are inexpensive to assess; models should abstain
from point estimation when support is insufficient.

\textbf{Limitations.} This study reanalyzed one biomarker for one disease pair. The 26
PVOD cases came from three registries, and no independent test set was available. We
therefore cannot estimate external generalizability, and the method has not been
validated for clinical use. The general contribution is a reporting approach for
individual support, not a clinically validated diagnostic rule.

\bibliographystyle{plain}
\begin{thebibliography}{1}
\bibitem{source}
A.~Prabhakar, E.~M.~D.~D. De~Bie, J.~T. DesJardin, P.~Ghatpande, S.~Gr\"af,
L.~S. Howard, S.~J. Wort, C.~Church, D.~G. Kiely, UK Pulmonary Arterial Hypertension
Cohort Study Consortium, E.~Sumpena, T.~Aung, S.~Carter, J.~Kukreja, S.~Hays,
J.~R. Greenland, J.~P. Singer, M.~Wax, P.~J. Wolters, M.~A. Simon, M.~Toshner,
G.~Lagna, and A.~Hata.
GDF15 is a putative biomarker for distinguishing pulmonary veno-occlusive disease and
pulmonary arterial hypertension.
\emph{J Clin Invest}, 136(2):e199159, 2026. Research Letter. doi:10.1172/JCI199159
\end{thebibliography}

\end{document}
